\section{Current new result and unchanged scope}
This section is the new contribution of this delivery. The preceding
424 body is retained in full, including its qualifications and historical
claims; retention is not a correctness grant. All four actual notes in
all five fields, their three distinct expansions, the separate 588-LF
proof, all 2390 LF of the standalone, current scope records, both reading
ledgers, verdict and residual have been read. All fifteen complete
immediate primary units have been freshly read, including their proofs.
The three supplied native programs were read as inert text only.
The entire deeper historical, primary, book, geometric and coherent-descent
audit is not complete. In particular the genuine nonalgebraization premise
is not newly certified. The supplied compact good strict smooth separated
boundaryless pure analytic quotient is input.

The original $\C$, $\mathcal G$, raw mass two, energy $E$ and full
dual $D$ remain unchanged. We prove a prescribed nonskeletal equation
sector, rather than merely insert a partial atom in an unspecified measure.
We also prove an integral representation result for dominated inputs.
The equation for $2\delta_x$ and the general arbitrary-support equation
remain open.

\section{Exact analytic inputs at the new applications}
CLD v2 Theorem 15.7.1 equips a topologically separated pure $n$-dimensional
analytic space, with its actual analytic boundary, with its calibrated
tropical calculus. Here $n=2$ and the ambient boundary is empty; all local
applications below are on ordinary open subsets, not compact domains with
their boundary suppressed. Definition 8.1.7 and Lemma 8.1.8 define the
entire minimal continuous sheaf by positive smooth generators and local
uniform closure. Corollary 8.2.4 permits convex coordinatewise increasing
composition, including finite maxima, of continuous psh slots.
Lemma 8.5.5 and 8.7.1 define actual positive recursive products;
8.5.1 and 8.7.2 give their local uniform continuity, 8.7.4 symmetry,
and 8.8.5 multilinearity and invariance under smooth $\DD$-null changes.
The fixed-current restriction formula is 8.5.6, iterated in 8.7.5:
for a topologically proper tropical morphism $j$ and positive closed $S$,
$\DD u\wedge j_*S=j_*(\DD(j^*u)\wedge S)$.
Its smooth integration-by-parts proof extends by uniform closure.
We use it for the closed fixed analytic fibres regarded over their
completed residue field, exactly as in CLD 18.6.1--2. No moving Gauss
section is asserted analytic. CLD 7.2.8--13 supplies positive top Radon
currents, coefficient measures and continuous-test density.

The decisive exact primary formula is CLD v2 Corollary 18.6.2:
on a boundaryless, topologically separated, equidimensional analytic
space, for analytic $f_1,\ldots,f_d$ and positive real
$t,r_1,\ldots,r_d$,
\[
 (\DD\log\max(t,|f_1|/r_1,\ldots,|f_d|/r_d))^d
 =\delta_{f^{-1}(\eta_{(tr_1,\ldots,tr_d)})}.
\]
The right side is the fixed-fibre integration current. Its proof first
removes points where the tropical dimension is too small, then uses
18.6.1 on the flat locus and computes the projective residual intersection
coefficient as one. The retained literal proof has typographic slips in
its first-case prose (the vanishing is that of the top power); we use
the displayed statement and its top-power argument, not those slips.
For our two-dimensional coordinate translations and inversions the fibre
is a singleton and the map is an isomorphism, of degree one. For $d=1$
the source is a surface and the fibre is an actual curve, and the final
application on that curve gives a singleton. All radii are positive;
strictness of a weighted-radius domain is not required by this formula.
Thus the coefficients below are one, not two. Equivalently the
identity moment has saturated primitive lattice and degree one by
14.5.2--3; a gradient triangle has area $1/2$ and the raw factor is $2!$.
The integration and reconstruction applications are 5.10.1(c),(d)
and 5.10.2. No extra extension-degree factor is inserted.

The retained bounded reconstruction uses smooth Green/Stokes 5.11.1
with zero global boundary term, finite coefficient measures and nested
cutoffs to include every continuous stage. Namely strong gradients
are closable since a uniformly zero sequence pairs to zero against every
smooth gradient test by closedness. For uniformly convergent bounded
psh functions $f_j$, with $|f_j|\le A$, cutoff $\chi$ and $\eta=1$
near its support, the strong-gradient estimate is
\[
 Q_{T,\chi}(f_j-f_k)\le \epsilon^2\|\DD\chi T\|_{\rm TV}/2
 +2A\epsilon\|\chi\|_\infty\|\DD\eta T\|_{\rm TV}.
\]
Consequently the gradient domain is uniformly closed and contains the
whole minimal sheaf. This gives $Q_T(f,g)=-\int f\DD(gT)\ge0$
on the diagonal, cubic differentiation, tangent concavity, and midpoint
gain $I/8$. The complete retained proof gives fixed bounded-test $L^1$
moduli on normalized energy sublevels and positive full-point atoms.
For clarity their ingredients are finite coordinate-curve paths with
collar caps, actual fixed-fibre currents, and Cauchy--Schwarz; full
points have Laurent Abhyankar coordinates, local singleton finite-flat
Gauss representatives (11.6.3/5(c),13.8.1--5), and a log-max cap.
The current transgression estimate is
\[
 Q_{\theta_w}(a-b)\le I(a,b)/2+\sqrt{42C I(a,b)}
\]
on a common nonpositive energy sublevel. The bound follows by writing
$g=(a+b)/2$, $h=w-g$ and bounding each transgression pairing's
$h$-gradient by $12C+3C+6C=21C$. A smooth absolute-value ramp,
Green and the finite-grid path estimate give the stated $L^1$ moduli;
an atom forces convergence at its full point. Laurent caps separate
actual points; the span of $\C$ is a vector lattice and the compact
finite maximum/minimum interpolation proof makes it dense in $C(X)$.
These are reproduced bounded deductions at the disclosed conventional
CLD/Berkovich/graded/flat/curve foundation boundary, not supplied verdicts.

For generated uniqueness and the ordinary upper relaxation we use GR v2
unnumbered Section 1.2 (CLD positivity implies classical curve-test
positivity), Definition 4.1 and Proposition 4.2 (all curves and extensions),
and Lemma 6.3/Theorem 6.4 (finite full-dimensional affinoid Shilov boundary
and its maximum). Their full immediate proofs were read. Smooth partition,
Berkovich affinoid bases, Temkin, Ducros, Wanner and Thuillier deeper
foundations remain expressly conventional. GR 6.5 is not used to create
a representative. CLD 8.6.7 is not applied.

\section{A global positive modification with genuine collar gluing}
Put $g=\sigma(0)$. Let $c>0$, $R=e^{-c}$, and let $x_c$ be the
joint Gauss point of $|z_i-1|\le R$. When $c=1/8$ this is exactly
the specified point $z_i=1+s u_i$. Its coordinate initials are independent;
it is a full point, $\tau(x_c)=0$, and $|z_i-1|(x_c)=R<1$
whereas $|z_i-1|(g)=1$.

Define the global invariant baseline
\[
 q_0=\tfrac12\xi_1(1-\xi_1)+\tfrac12\xi_2(1-\xi_2)-\tfrac14,
 \qquad \xi_i=r_i-\lfloor r_i\rfloor.
\]
The convex piecewise-affine interpolant of $r_i^2/2$ at integers has
positive slope jumps one. Positive periodic convolution gives actual
smooth positive weights converging uniformly to $q_0$, so $q_0\in\C$.
In the injective chart $U=\{|r_i|<1/4\}$,
\[
 \rho+q_0=f_1+f_2-\tfrac14,
 \quad f_i=\log A_i-\tfrac12\log|z_i|,
 \quad A_i=\max(1,|z_i|)=\max(1,|z_i-1|).
\]
Write $w_i=z_i-1$, and set on $U$
\[
 L_c=\log\max(R,|w_1|,|w_2|,|w_1w_2|),\qquad
 v_c=L_c-\log(A_1A_2).
\]
Extend $v_c$ by zero outside $U$ and put $q_c=q_0+v_c$.

\begin{lemma}\label{newns:glue}
This is an actual continuous global weight in the entire original $\C$.
One has $-c\le v_c\le0$, $v_c(g)=0$, $v_c(x_c)=-c$;
$v_c=0$ unless both $|w_i|<1$, in particular unless $\tau=0$.
It vanishes on an ordinary open collar of the chart boundary.
\end{lemma}
\begin{proof}
If $|w_1|\ge1$, the maximum defining $L_c$ is
$|w_1|\max(1,|w_2|)=A_1A_2$; the other coordinate is identical.
If both are less than one, then $A_i=1$ and
$v_c=\log\max(R,|w_1|,|w_2|)\in[-c,0]$.
If $|z_i|\ne1$, ultrametricity gives either $|w_i|=|z_i|>1$
or $|w_i|=1$. Thus the correction is zero there.
On the ordinary collar $1/8<\max|r_i|<1/4$ it is identically zero;
this, rather than an assertion about Shilov vertices, proves gluing.
The local new potential is
\[
 \rho+q_c=L_c-\tfrac12\log|z_1z_2|-\tfrac14.
\]
The constant $R>0$ keeps the logarithmic maximum finite and continuous.
Near a zero of a polynomial slot, that slot is strictly below the
positive constant and is omitted on a neighborhood. The remaining
nonzero slots are unit logarithms; smooth convex maxima approximate
$L_c$ locally uniformly by ordinary positive CLD generators.
The subtracted unit logarithms are smooth and $\DD$-null. Minimal
continuous-sheaf closure proves positivity at every point, including
polynomial zeros. On the collar the new potential equals the old one,
so sheafiness glues. Deck units $2,3$ do not require descent of $z_i-1$:
that function is used only on this injective chart and the correction
already vanishes on the collar.
\end{proof}

\section{Exact analytic measures and a prescribed nonskeletal equation}
\begin{lemma}\label{newns:products}
As actual measures on $X$, on all continuous analytic tests,
\[
 \MA(q_0)=2\delta_g,\qquad
 \MA(q_c)=\delta_g+\delta_{x_c},\qquad
 B(q_0,q_c)=2\delta_g.
\]
\end{lemma}
\begin{proof}
Locally $\DD f_i=\DD\log\max(1,|z_i|)$ is the integration
current of the fixed coordinate Gauss fibre by 18.6.2 with $d=1$.
Its self-square is zero by the one-coordinate dimension rule, and the
product of the two fibres is the singleton joint Gauss point with
coefficient one (apply the same formula on the first fibre).
This gives the baseline square $2\delta_g$. Between integer coordinate
hyperplanes the baseline local potential is affine; elsewhere a single
coordinate curvature has zero square. This covers the quotient and
proves there is no other baseline mass.

On $U$, the harmonic change reduces the new square to $(\DD L_c)^2$.
Where some $w_i=0$, locally discard its two slots and use only one
coordinate; the square is zero. Otherwise put $b_i=\log|w_i|$.
The affine maximum $\max(-c,b_1,b_2,b_1+b_2)$ has at most two
active slots except at the two vertices $(-c,-c)$ and $(0,0)$.
Indeed the triples containing the constant and the product along with
one single coordinate are dominated by the other single coordinate.
With at most two active slots, the maximum locally is a unit logarithm
plus a one-variable maximum. Smooth approximation and dimension one
make its square zero. At $(-c,-c)$ the product slot is strictly smaller;
locally $L_c=\log\max(R,|w_1|,|w_2|)$.
Apply 18.6.2 with $d=2$, $t=R$, $r_i=1$, $f_i=w_i$:
the fibre is exactly the singleton $x_c$, degree one.
At $(0,0)$ the constant slot is strictly smaller; subtracting
the harmonic $\log|w_1w_2|$ leaves
$\log\max(1,|w_1|^{-1},|w_2|^{-1})$.
Apply the same theorem with $t=r_i=1$, $f_i=w_i^{-1}$.
Its singleton Gauss fibre is $g$; translation by one preserves the
unit bidisc Gauss norm. These identities also remove all non-Gauss
points with the same logarithmic radii. No radial projection is used
to define a measure. Outside $U$ the potential is unchanged and the
baseline measure is zero. Hence the second displayed identity holds
globally, with both coefficients one and total mass two.

For the mixed product restrict $L_c$ to the fixed fibre
$z_1=\eta_1$ over $\mathscr H(\eta_1)$. The element $z_1-1$ of
this completed field has norm one at every point of that fibre, so
$L_c=\log\max(1,|w_2|)=\log\max(1,|z_2|)$ there.
The harmonic correction is unchanged. The fixed-current restriction
formula therefore gives $\DD f_1\wedge\theta_{q_c}=\delta_g$.
The reversed coordinate gives the other $\delta_g$.
Outside $U$, $B(q_0,q_c)=\MA(q_0)=0$.
These are equalities of recursive currents, and 7.2.8--13 make them
equalities of Radon measures against every continuous test.
\end{proof}

\begin{theorem}\label{newns:equation}
For every $c>0$ and every prescribed $0\le m\le1$, the nonskeletal
target (when $m>0$)
\[
 \mu_{m,c}=(2-m)\delta_g+m\delta_{x_c}
\]
has the actual normalized bounded solution
$q_{m,c}=q_0+\sqrt m\,v_c\in\C\subset\mathcal G$.
It is unique modulo constants among all actual $\mathcal G$ solutions.
It attains the full dual, with
\[
 E(q_{m,c})=-\tfrac13-\tfrac c3m^{3/2},\qquad
 D(\mu_{m,c})=\tfrac16+\tfrac{2c}{3}m^{3/2}.
\]
\end{theorem}
\begin{proof}
Let $\lambda=\sqrt m\in[0,1]$. Convexity of the actual positivity
class and Lemma \ref{newns:glue} give $q_{m,c}\in\C$.
Multilinearity and Lemma \ref{newns:products} give
\[
 \MA((1-\lambda)q_0+\lambda q_c)
 =(2-\lambda^2)\delta_g+\lambda^2\delta_{x_c}=\mu_{m,c}.
\]
Since the correction is nonpositive and vanishes at the baseline's
maximum, $\sup q_{m,c}=0$; its range is in $[-1/4-c\sqrt m,0]$.
An actual decreasing presentation is $q_{m,c}+2^{-n}\downarrow q_{m,c}$,
with finite lower energy.

For the baseline, the reference, mixed-reference and top moments are
$-1/6,-1/3,-1/2$. The one-coordinate mean is
$\int_0^1[\xi(1-\xi)/2-1/8]d\xi=-1/24$;
the reference is $2\sigma_*dr_1dr_2$, the mixed-reference is the sum
of the two coordinate-circle Lebesgue measures, and the top is $2\delta_g$.
These measures follow from the same fixed-fibre calibration and smooth
reference integration. Thus $E(q_0)=-1/3$.
Along the actual segment, first variation is
$\frac d{d\lambda}E(q_0+\lambda v_c)=\int v_c\MA(q_0+\lambda v_c)
=-c\lambda^2$. Integration gives the stated energy.
The target moment is $-1/2-cm^{3/2}$, since the two target values are
$-1/4$ and $-1/4-c\sqrt m$.
Tangent concavity on the entire original $\C$ gives
$E(a)\le E(q_{m,c})+\int(a-q_{m,c})\mu_{m,c}$ for every $a$.
Thus the full dual, without an invariant-test restriction, is attained
and equals the displayed finite value.

Here is generated uniqueness with its actual representation antecedent
retained. For $c_n\downarrow v$ in $\mathcal G$, bounded lower energy,
compact strict domination compares any fixed term of one presentation
plus $\epsilon$ with a sufficiently late term of another presentation.
Energy monotonicity and $E(a+\epsilon)=E(a)+2\epsilon$ make the
limiting energy canonical. Apply this also to finite convex combinations
and evaluate their degree-three energy polynomials at a finite unisolvent
set to get canonical coefficient limits. For a fixed bounded $w$,
$P_n(t)=E((1-t)c_n+tw)$ and $A_n(t)=E(tc_n)$ give
$\int w\MA(c_n)=P'_n(0)+A'_n(1)$.
Dense tests and common positive mass give the canonical weak operator
$\MA(v)=\lambda$ and $d=\lim\int c_n\MA(c_n)$.
The bounded tangent inequality yields $D(\lambda)\le e-d$.
Monotone convergence in its dual inequality gives $v\in L^1(\lambda)$.
The ordered cocycle gives
$0\le\int(c_n-c_l)\MA(c_l)\le3(E(c_n)-E(c_l))$, $l\ge n$.
Pass $l$ then $n$ to infinity: $d=\int v\lambda$, so the presentation
maximizes $D(\lambda)$.
Its suprema decrease to the finite $s=\sup v$ by nested compact
superlevel sets and $E(c_n)\le2\sup c_n$.
The normalized $a_n=c_n-\sup c_n$ converge everywhere to $v-s$.
If $\lambda=\MA(q_{m,c})$, midpoint gain and full-point moduli force
$v-s=q_{m,c}$ at every full point. Such points are dense by the
finite Shilov boundary in every pure affinoid neighborhood.
USC of the already decreasing limit prevents $v-s<q_{m,c}$ anywhere.
For the opposite inequality pass $n$ to infinity in the maximum bound
on each finite affinoid Shilov boundary, then shrink the affinoid
neighborhood using continuity of $\rho+q_{m,c}$. This gives
$v-s\le q_{m,c}$ everywhere. Hence uniqueness holds on all of $X$.
\end{proof}

\begin{corollary}\label{newns:nonordered}
For every normalized nonordered maximizing net for $D(\mu_{m,c})$,
the canonical full-point tuple has the actual continuous representative
$q_{m,c}$, the canonical ordinary upper relaxation equals this weight,
and the genuine bounded MA measures tend to $\mu_{m,c}$ on every
continuous analytic test. Every fast sequence has $J=K=0$.
This is representation for this sector, not a general recognition theorem.
\end{corollary}
\begin{proof}
Coercivity gives a common energy sublevel. Midpoint gain with the
attained competitor gives $I(a_i,q_{m,c})\le4\varepsilon_i$.
The reconstructed moduli imply energy, fixed-test $L^1$, full-point
and weak-MA convergence. Every finite affinoid Shilov boundary consists
of full points. The maximum bound followed by shrinking neighborhoods
shows that the ordinary upper relaxed net limit is at most the continuous
$q_{m,c}$; conversely every such neighborhood contains a full point
whose limiting value is arbitrarily close to $q_{m,c}(x)$.
Thus the relaxation is exactly the already constructed actual weight.
Fast Borel limits agree at the two full target atoms, hence in $L^1$
of the target; the limiting signed moment is the target moment of the
solution and the weak measure is the target, proving $J=K=0$.
We do not assign positivity to the raw Borel limit at other points.
\end{proof}

\section{An exact obstruction to extending this segment to the whole Dirac}
\begin{proposition}\label{newns:obstruction}
For $\lambda>1$, the proposed extension $q_0+\lambda v_c$ is not
in $\C$. In particular this segment cannot reach the whole-mass target
$2\delta_{x_c}$ (its formal square would require $\lambda=\sqrt2$).
This is an obstruction to this explicit segment, not a nonexistence theorem.
\end{proposition}
\begin{proof}
Choose $a\in K$ with $0<|a|<R$ and restrict to the actual analytic
curve $w_2=a$. Near the unit Gauss point of $w_1$, $A_2=1$ and
$L_c=\log|w_1|$: the constant and the $w_2$ and product slots are
strictly smaller. The local potential on this curve, up to smooth
harmonic terms and a constant, is
\[
 (1-\lambda)\log\max(1,|w_1|)+\lambda\log|w_1|.
\]
Its $\DD$ is $(1-\lambda)\delta_{\eta_1}$ by the one-dimensional
fixed-fibre formula. Both $w_1$ and $z_1=1+w_1$ are nonzero near
this Gauss point, so the discarded logarithms really are smooth harmonic.
The negative atom contradicts positivity of the pullback of a member of
the minimal continuous class (positive smooth generators pull back and
uniform closure preserves positivity). The coordinate curve lies in the
injective chart near this point and has no ambient boundary there.
\end{proof}

\section{Domination gives prescribed-measure absolute convergence}
\begin{theorem}\label{newns:domination}
Let $\mu$ be any positive raw-mass-two Radon measure with
$\mu\le A\MA(w)$, where $w\in\C$ is fixed bounded and $A<\infty$.
Then $A\ge1$ and, writing $B_w=D(\MA(w))$,
\[
 D(\mu)\le A(1+B_w)+\tfrac32A^2B_w<\infty.
\]
Every normalized maximizing family is Cauchy in $L^1(\mu)$ with a
canonical limit. For every fast sequence with the earlier increments
$\int|a_{n+1}-a_n|\MA(w)\le2^{-4n}$ for all sufficiently late $n$,
its actual Borel $u$ satisfies, for all sufficiently late $j$,
\[
 \int|a_j-u|\,d\mu\le A\sum_{n\ge j}2^{-4n}
 =\tfrac{16A}{15}2^{-4j},\qquad J=0.
\]
It annihilates the exceptional $N$ and its negative sequence is uniformly
integrable. Neither $u\in\mathcal G$ nor $\nu=\mu$ follows here.
\end{theorem}
\begin{proof}
Normalize $w$. For normalized $a$ put $C=-E(a)$,
$A_0=\int(-a)\theta^2\le1$, $q_0=Q_\theta(a)$ and
$q_1=Q_{\theta_a}(a)$. Cubic differentiation gives
$C=A_0+2q_0/3+q_1/3$ and
$-E(ta)=tA_0+t^2(1-t/3)q_0+t^3q_1/3\le t+3t^2C/2$.
Apply the full dual of $\MA(w)$ to $ta$, divide by $t$, and choose
$t=\min(1,\sqrt{2B_w/(3C)})$, with the limiting zero cases.
Then $\int(-a)\MA(w)\le1+B_w+\sqrt{6B_wC}$.
Domination bounds $F_\mu(a)$ by
$-C+A(1+B_w)+A\sqrt{6B_wC}$; maximizing the scalar quadratic
in $\sqrt C$ proves the dual bound. Equal masses give $A\ge1$.

On a common maximizing energy sublevel, midpoint gives
$I(a_i,a_j)\le4(\varepsilon_i+\varepsilon_j)$.
The fixed-$w$ absolute modulus and domination give $L^1(\mu)$
Cauchyness. Completeness and comparison of two families on product
indices give a canonical limit without countable cofinality.
For a fast sequence, choose an integer $k$ with normalized $w$ in
$\C_{C+k}$. Its exceptional capacity bound implies
$\MA(w)(N)=0$, hence $\mu(N)=0$.
Countable Tonelli gives
$\int\sum_{n\ge j}|a_{n+1}-a_n|\,d\mu\le A\sum_{n\ge j}2^{-4n}$.
Outside $N$ the actual pointwise limit is $u$, and
$|a_j-u|$ is bounded by that sum. This proves the displayed actual
$L^1$ bound; canonicality identifies this limit with the one for every
maximizing family. Its integral is $e-D(\mu)$, so $J=0$.
Uniform integrability for a sequence follows directly from $L^1$
convergence: $(v_n-T)_+\le|v_n-v|+(v-T)_+$, followed by control
of a late tail and then finitely many initial integrable functions.
No inference from null-set annihilation alone was made.
\end{proof}

In particular, for the unchanged $x=x_{1/8}$, the new explicit equation
$\MA(q_{1,1/8})=\delta_g+\delta_x$ gives
$2\delta_x\le2\MA(q_{1,1/8})$.
Thus every fast family for the originally proposed whole-mass Dirac
has $J=0$ and the quantitative integral bound with $A=2$.
The exact dual of the dominating measure is $1/4$, yielding
\[
 \tfrac38\le D(2\delta_x)\le4.
\]
The lower bound evaluates $q_{1,1/8}$: its energy is $-3/8$ and
its value at $x$ is $-3/8$. The upper bound is the domination bound.
We have not computed this Dirac dual exactly, proved $K=0$, or produced
an actual solution for it. Its general nonordered analytic representation
gate remains open despite its now-proved vanishing integral defect.

\section{Current verdict}
The explicit nonskeletal two-atom sector has an actual original-$\C$
equation, actual decreasing presentation, full-dual attainment, energy,
generated uniqueness and representation of the canonical nonordered
maximizing tuple. Its analytic bridges are the finite log-max generators,
ordinary collar gluing, exact fixed-fibre products and calibrated raw mass.
It is narrower than unchanged P16.7. Domination solves the prescribed
$L^1$ loss issue, including for $2\delta_x$, but does not solve its
analytic representation or equation. Adequate general envelopes,
arbitrary-support existence and genuine nonalgebraization/descent remain
open. The separate original residual is byte-unchanged. All material is
an unverified local draft; no compiler, publication or acceptance is claimed.
Candidate is false.
